\chapter{Parent-origin источника и жёсткости центральной щели}
\label{ch:v8-baryon-c0-singlet-triplet-central-gap-source-stiffness-parent-origin-gate}

Минимальная архитектура требует двух коэффициентов $j$ и $m^2$.
Проверим отдельно существование нужных следовых форм и происхождение их
численных коэффициентов из уже построенного родителя.

\section{Фактическое ограничение старого родителя}

Центральное направление удовлетворяет
\begin{equation}
 Q=P_{\mathbf3}-\frac34I_4,
 \qquad \Tr Q=0,
 \qquad \Tr Q^2=\frac34.
 \label{eq:v8-baryon-c0-gap-origin-q-invariants}
\end{equation}
Но сама семейная тройка и её ковариантный endpoint были лишь условными
расширениями:
\begin{equation}
 N_{SO(3)\,\rm origin}=0/5,
 \qquad N_{\rm endpoint\ extension}=0/3.
 \label{eq:v8-baryon-c0-gap-origin-missing-endpoint-ledgers}
\end{equation}
Поэтому старое действие не зависит от новой центральной координаты и его
нулевое продолжение даёт
\begin{equation}
 \left.\frac{\partial S_{\rm old}}{\partial\lambda}\right|_0=0,
 \qquad
 \left.\frac{\partial^2S_{\rm old}}{\partial\lambda^2}\right|_0=0,
 \qquad (j,m^2)_{\rm old}=(0,0).
 \label{eq:v8-baryon-c0-gap-origin-old-parent-zero-restriction}
\end{equation}

\section{Незавешенный спектральный след}

Пусть $H(\lambda)=\varepsilon I_4+\lambda Q$. Для каждого целого $n\geq1$
линейный отклик обычного момента равен
\begin{equation}
 \left.\frac{d}{d\lambda}\Tr H(\lambda)^n\right|_{\lambda=0}
 =n\varepsilon^{n-1}\Tr Q=0.
 \label{eq:v8-baryon-c0-gap-origin-unmarked-moment-no-source}
\end{equation}
Следовательно, никакой аналитический незавешенный след вокруг скалярного
фона не создаёт $j$. Второй отклик имеет вид
\begin{equation}
 \left.\frac{d^2}{d\lambda^2}\Tr H(\lambda)^n\right|_{\lambda=0}
 =\frac34n(n-1)\varepsilon^{n-2}.
 \label{eq:v8-baryon-c0-gap-origin-unmarked-moment-curvature}
\end{equation}
В частности, в начале координат для
$S=\sum_{n\geq1}a_n\Tr H^n$ остаётся
\begin{equation}
 \left.S''\right|_{\varepsilon=\lambda=0}=\frac32a_2.
 \label{eq:v8-baryon-c0-gap-origin-quadratic-trace-stiffness}
\end{equation}
След фиксирует форму нормы, но не размерный коэффициент $a_2$.

Линейный обычный след исчезает непосредственно:
\begin{equation}
 \Tr(\lambda Q)=0.
 \label{eq:v8-baryon-c0-gap-origin-ordinary-linear-trace-zero}
\end{equation}

\section{Маркированные следы как условный near miss}

Центрированные grading и семейный оператор Казимира совпадают,
\begin{equation}
 \Gamma_4^0=\mathcal C_2^0=2Q,
 \qquad
 \Tr(\Gamma_4^0\lambda Q)
 =\Tr(\mathcal C_2^0\lambda Q)=\frac32\lambda.
 \label{eq:v8-baryon-c0-gap-origin-marked-linear-responses}
\end{equation}
Поэтому условный функционал
\begin{equation}
 S_{a_2,g}[h^0]=a_2\Tr((h^0)^2)-g\Tr(\Gamma_4^0h^0)
 \label{eq:v8-baryon-c0-gap-origin-conditional-marked-parent}
\end{equation}
действительно реализует нужную форму. При $h^0=\lambda Q$ он даёт
\begin{equation}
 S_{a_2,g}=\frac34a_2\lambda^2-\frac32g\lambda,
 \qquad m^2=\frac32a_2,
 \qquad j=\frac32g.
 \label{eq:v8-baryon-c0-gap-origin-conditional-coefficient-map}
\end{equation}
Однако $\Gamma_4^0$ и $\mathcal C_2^0$ задают одну метку, а не два
независимых источника, и коэффициенты $a_2,g$ отсутствуют в текущем
действии. Два точных свидетеля сохраняют один вакуум:
\begin{equation}
 (a_2,g)=\left(\frac23,\frac23\right)\Longrightarrow(m^2,j)=(1,1),
 \qquad
 (a_2,g)=\left(\frac43,\frac43\right)\Longrightarrow(m^2,j)=(2,2).
 \label{eq:v8-baryon-c0-gap-origin-free-coefficient-witnesses}
\end{equation}

Нескалярный фон мог бы породить линейный отклик, но только ценой уже
введённой щели. Например,
\begin{equation}
 B=\varepsilon I_4+bQ
 \quad\Longrightarrow\quad
 \left.\frac{d}{d\lambda}\Tr(B+\lambda Q)^2\right|_0
 =\frac32b.
 \label{eq:v8-baryon-c0-gap-origin-circular-background-source}
\end{equation}
Использовать $b\ne0$ для происхождения $j$ означало бы предположить искомый
центральный фон заранее.

\section{Типовые и тепловые препятствия}

Полный шумовой гессиан живёт на старом $42$-мерном носителе, тогда как
новая координата принадлежит $\mathbb RQ$. Каноническое отображение между
ними не построено:
\begin{equation}
 K_{\rm noise}\in\operatorname{End}(\mathbb R^{42}),
 \qquad h^0\in\mathbb RQ,
 \qquad \iota:\mathbb RQ\longrightarrow\mathbb R^{42} \text{не задано}.
 \label{eq:v8-baryon-c0-gap-origin-no-noise-hessian-pullback}
\end{equation}
KMS-условие также действует после выбора гамильтониана:
\begin{equation}
 \frac{\gamma_+}{\gamma_-}=e^{-\beta\lambda_*},
 \qquad
 \lambda_*=\frac{j}{m^2},
 \label{eq:v8-baryon-c0-gap-origin-kms-postselection}
\end{equation}
и потому не определяет $j$ и $m^2$ как коэффициенты действия.

Итоговое разделение равно
\begin{equation}
 N_{\rm conditional\ shapes}=2/2,
 \qquad
 N_{\rm parent\ origin}=0/8.
 \label{eq:v8-baryon-c0-gap-origin-final-ledgers}
\end{equation}

\begin{theorem}[Следовой запрет источника на скалярном фоне]
Для любого аналитического незавешенного следового функционала от
$\varepsilon I_4+\lambda Q$ первая вариация по $\lambda$ в нуле равна
нулю. Квадратичный след условно предоставляет форму жёсткости, а
grading/Casimir-вставка --- форму источника, но их коэффициенты и сама
новая endpoint-структура не принадлежат текущему родителю.
\end{theorem}

\begin{nogocriterion}[Запрет кругового фонового источника]
Нельзя выводить $j$ расширением следа вокруг фона с ненулевой компонентой
$bQ$: такая компонента уже является искомой центральной щелью. Аналогично,
нельзя переносить старый шумовой гессиан без типизированной карты нового
endpoint в $42$-мерный носитель.
\end{nogocriterion}

\begin{protocol}\begin{enumerate}
\item фактические $(j,m^2)$ старого родителя: \textbf{$(0,0)$};
\item незавешенный аналитический след создаёт $j$: \textbf{нет};
\item квадратичный след создаёт форму жёсткости: \textbf{да условно};
\item grading/Casimir-вставка создаёт форму источника: \textbf{да условно};
\item независимых маркированных направлений: \textbf{$1$};
\item канонический pullback шумового гессиана: \textbf{нет};
\item условных форм: \textbf{$2/2$};
\item выведенных каналов происхождения: \textbf{$0/8$};
\item следующий гейт: \textbf{допуск динамического носителя источника}.
\end{enumerate}\end{protocol}

Машинный сертификат сохранён в
\path{s2t_v8_baryon_c0_singlet_triplet_central_gap_source_stiffness_parent_origin_gate_results.json}.